% =====================================================================
%  CHƯƠNG 2 — KIẾN TRÚC KNOWLEDGE GRAPH VÀ CƠ SỞ TRI THỨC
% =====================================================================
\chapter{Kiến trúc Knowledge Graph và cơ sở tri thức}\label{chap:kg}

Chương~\ref{chap:mo-dau} định nghĩa KG là một tập bộ ba. Chương này mở tập bộ ba đó ra: một KG được cấu tạo từ
những gì, được viết ra theo chuẩn RDF thế nào, bản thể học thêm vào đó những ràng buộc và luật suy diễn ra sao,
và khi nào nên lưu tri thức trong một cơ sở dữ liệu đồ thị thay vì một cơ sở dữ liệu vector. Mọi ví dụ đều lấy
từ hai KG thật của bài, cài trong \texttt{tl2/kg\_base.py}.

\section{Các thành phần cấu tạo của Knowledge Graph}
Một KG có ba loại thành phần (Hình~\ref{fig:thanhphan}):

\begin{itemize}
  \item \textbf{Thực thể (entity)} là nút của đồ thị, đại diện cho một đối tượng có danh tính: một bệnh
        (\textsf{DTD\_Tip2}), một thuốc (\textsf{Metformin}), một món ăn (\textsf{Com\_Trang}), một chính sách
        (\textsf{ND\_232\_2025}), một cổ phiếu (\textsf{FPT}). Mỗi thực thể có mã định danh duy nhất, một
        \textbf{kiểu} là một lớp trong bản thể học, nhãn tiếng Việt để hiển thị và danh sách bí danh để liên kết
        với câu hỏi tự nhiên (``tiểu đường'', ``đtđ'', ``type 2'' đều trỏ về \textsf{DTD\_Tip2}).
  \item \textbf{Quan hệ (relation)} là cạnh có hướng, gán nhãn bằng một vị từ:
        $\textsf{Metformin} \xrightarrow{\textsf{TUONG\_TAC\_NGUY\_HIEM}} \textsf{Ruou\_Bia}$. Giữa hai thực thể có
        thể có nhiều cạnh với vị từ khác nhau, nên cấu trúc dữ liệu phù hợp là \emph{đa đồ thị có hướng}
        (\texttt{networkx.MultiDiGraph}~\cite{hagberg2008networkx}).
  \item \textbf{Thuộc tính (attribute)} gắn lên nút (mã ICD-10 \texttt{E11}, chỉ số GI 73 của cơm trắng, ngưỡng
        126 mg/dL) hoặc lên cạnh (mức độ \texttt{cao}, hậu quả, nguồn trích dẫn \texttt{h03-ruou-bia}).
\end{itemize}

\hinhtikz{t2_components}{Ba thành phần của một KG, minh họa trên \texttt{VietnamDiabetesKG}: thực thể là nút
có kiểu và thuộc tính, quan hệ là cạnh có nhãn vị từ, thuộc tính có thể nằm trên cả nút lẫn cạnh}{thanhphan}

Thuộc tính trên cạnh là chỗ hai mô hình dữ liệu đồ thị phổ biến khác nhau. \textbf{Đồ thị thuộc tính có nhãn}
(Labeled Property Graph, LPG), mô hình của Neo4j, cho phép gắn trực tiếp cặp khóa--giá trị lên cạnh.
\textbf{RDF} chỉ có bộ ba, nên muốn nói ``cạnh này có mức độ cao'' phải biến chính cạnh thành một thực thể, kỹ
thuật gọi là \emph{reification} (Hình~\ref{fig:lpgrdf}). Lớp \texttt{KnowledgeGraph} của bài làm việc theo LPG
khi xây dựng và truy vấn, rồi xuất sang RDF có reification khi cần suy diễn và SPARQL; nhờ vậy dùng được thế mạnh
của cả hai.

\hinhtikz{t2_lpg_rdf}{Cùng một sự kiện ``metformin tương tác nguy hiểm với rượu bia, mức độ cao'' trong mô hình
LPG (thuộc tính nằm trên cạnh) và mô hình RDF (cạnh được reification thành một \texttt{rdf:Statement})}{lpgrdf}

Bản thể học được cài bằng hai lớp dữ liệu nhỏ. \texttt{Predicate} mô tả một vị từ: miền (domain) là lớp của chủ
thể, miền giá trị (range) là lớp của đối tượng, cùng ba tính chất logic đối xứng, bắc cầu, nghịch đảo.
\texttt{Ontology} giữ cây lớp dưới dạng từ điển \emph{lớp $\to$ lớp cha} và trả lời câu hỏi ``lớp $A$ có phải
là một loại $B$ không'' bằng cách đi ngược lên gốc (Đoạn mã~\ref{code:kg_base__Ontology}).

\maNguon{kg_base__Predicate}{Lớp \texttt{Predicate}: vị từ kèm miền, miền giá trị và tính chất logic}
\maNguon{kg_base__Ontology}{Lớp \texttt{Ontology}: cây phân cấp lớp và phép kiểm tra \texttt{is\_a}}

Điểm quan trọng nhất của lõi KG là mọi bộ ba đều được \textbf{kiểm tra ràng buộc trước khi nạp}
(Đoạn mã~\ref{code:kg_base__KnowledgeGraph.validate}): chủ thể phải thuộc miền của vị từ, đối tượng phải thuộc
miền giá trị, xét theo cả lớp con. Thêm $(\textsf{Rau\_Xanh}, \textsf{CHI\_DINH\_DAU\_TAY}, \textsf{Metformin})$
bị từ chối ngay vì rau xanh không phải một bệnh. Mục~\ref{sec:rangbuoc} thử bốn bộ ba cố tình sai như vậy.

\maNguon{kg_base__KnowledgeGraph.validate}{Kiểm tra ràng buộc miền và miền giá trị khi thêm một bộ ba}

Để LLM đọc được tri thức, mỗi bộ ba được \textbf{diễn đạt thành câu} tiếng Việt kèm các thuộc tính cạnh
(Đoạn mã~\ref{code:kg_base__KnowledgeGraph.verbalize}). Đây là định dạng mà tác tử đưa vào prompt ở
Chương~\ref{chap:agent}, ví dụ ``Metformin —[tương tác nguy hiểm với]→ Rượu bia (ethanol) (severity: cao; effect:
tăng nguy cơ nhiễm toan lactic\dots)''.

\maNguon{kg_base__KnowledgeGraph.verbalize}{Diễn đạt bộ ba thành câu tiếng Việt để đưa vào ngữ cảnh LLM}

\section{Mô hình biểu diễn bộ ba RDF}
RDF (Resource Description Framework) là chuẩn W3C để mô tả tài nguyên trên Web bằng bộ
ba~\cite{w3c2014rdf}. Mỗi thành phần của bộ ba là một \textbf{IRI} (định danh tài nguyên toàn cục), một
\textbf{literal} (giá trị có kiểu như số, chuỗi có nhãn ngôn ngữ) hoặc một \textbf{nút trống}. Chủ thể là IRI hoặc
nút trống, vị từ luôn là IRI, đối tượng là bất kỳ loại nào. Một đồ thị RDF đơn giản là \emph{tập hợp} các bộ ba
(Hình~\ref{fig:bobardf}); hai đồ thị gộp lại bằng phép hợp tập hợp, điều làm RDF đặc biệt hợp để tích hợp dữ liệu
từ nhiều nguồn.

\hinhtikz{t2_triple}{Bộ ba RDF (chủ thể, vị từ, đối tượng) viết bằng IRI đầy đủ trong không gian tên
\texttt{ex:} của bài, và đồ thị RDF như một tập các bộ ba}{bobardf}

Phép ánh xạ từ KG của bài sang RDF (Đoạn mã~\ref{code:kg_base__KnowledgeGraph.to_rdf}) đi theo bốn bước:

\begin{enumerate}
  \item Mỗi lớp thành một \texttt{rdfs:Class}, quan hệ lớp cha thành \texttt{rdfs:subClassOf}.
  \item Mỗi vị từ thành một \texttt{rdf:Property} kèm \texttt{rdfs:domain}, \texttt{rdfs:range}, nhãn tiếng Việt và
        các tiên đề OWL \texttt{owl:SymmetricProperty}, \texttt{owl:TransitiveProperty}, \texttt{owl:inverseOf}.
  \item Mỗi thực thể có bộ ba \texttt{rdf:type} trỏ tới lớp, \texttt{rdfs:label} và một bộ ba cho mỗi thuộc tính nút.
  \item Mỗi cạnh thành một bộ ba; nếu cạnh có thuộc tính thì thêm một \texttt{rdf:Statement} mang
        \texttt{rdf:subject}, \texttt{rdf:predicate}, \texttt{rdf:object} và từng thuộc tính.
\end{enumerate}

\maDai{kg_base__KnowledgeGraph.to_rdf}{Xuất KG sang RDF bằng rdflib, giữ thuộc tính cạnh bằng reification}

Bảng~\ref{tab:kgstats} tổng hợp kích thước hai KG. Số bộ ba RDF lớn hơn nhiều số cạnh vì mỗi thực thể sinh thêm
bộ ba kiểu và nhãn, mỗi cạnh có thuộc tính sinh thêm khối reification.

\bang{tab_kg_stats}{Quy mô hai đồ thị tri thức của tiểu luận}{kgstats}

\section{Bản thể học và biểu diễn tri thức hình thức}
Theo định nghĩa kinh điển của Gruber, bản thể học là \emph{``một đặc tả tường minh của một khái niệm
hóa''}~\cite{gruber1993ontology}: nó liệt kê các loại sự vật tồn tại trong một miền và các quan hệ được phép giữa
chúng. Trong thực hành Web ngữ nghĩa, bản thể học được viết bằng RDFS và OWL~\cite{w3c2012owl} và gồm ba lớp
tri thức:

\begin{itemize}
  \item \textbf{Phân cấp lớp}: \textsf{SafeFood} $\sqsubseteq$ \textsf{Food} $\sqsubseteq$ \textsf{Substance}
        $\sqsubseteq$ \textsf{MedicalEntity}. Hình~\ref{fig:ontology} là cây lớp của \texttt{VietnamDiabetesKG}
        với đúng \SL{kg/health/classes} lớp khai báo trong mã.
  \item \textbf{Ràng buộc vị từ}: \textsf{TUONG\_TAC\_NGUY\_HIEM} có miền \textsf{Drug} và miền giá trị
        \textsf{RiskSubstance}; \textsf{KHUYEN\_NGHI\_AN} nối \textsf{Disease} với \textsf{Food}.
  \item \textbf{Tiên đề tính chất}: vị từ có thể đối xứng, bắc cầu hoặc là nghịch đảo của vị từ khác
        (Hình~\ref{fig:tiende}).
\end{itemize}

\hinhtikz{t2_ontology}{Cây phân cấp lớp (\texttt{rdfs:subClassOf}) của bản thể học y tế: đúng
\SL{kg/health/classes} lớp khai báo trong \texttt{health\_ontology()}}{ontology}

\hinhtikz{t2_axioms}{Ba tiên đề tính chất của vị từ: bắc cầu, đối xứng, nghịch đảo; cạnh suy ra vẽ nét đứt}{tiende}

\subsection{Suy diễn tiến trên RDF}
Bản thể học không chỉ để kiểm tra mà còn để \textbf{suy ra} tri thức không được ghi tường minh. Hàm
\texttt{rdfs\_closure} cài suy diễn tiến (forward chaining): áp lặp lại năm luật cho tới khi không sinh thêm bộ ba
mới nào (điểm bất động). Với $\sqsubseteq$ là \texttt{rdfs:subClassOf} và \texttt{a} là \texttt{rdf:type}:
\begin{align}
  &(A \sqsubseteq B) \wedge (B \sqsubseteq C) \;\Rightarrow\; A \sqsubseteq C &&\text{(phân cấp bắc cầu)}\\
  &(x\ \texttt{a}\ A) \wedge (A \sqsubseteq B) \;\Rightarrow\; x\ \texttt{a}\ B &&\text{(kế thừa kiểu)}\\
  &p \text{ đối xứng} \wedge (x\, p\, y) \;\Rightarrow\; (y\, p\, x) &&\text{(đối xứng)}\\
  &p \text{ bắc cầu} \wedge (x\, p\, y) \wedge (y\, p\, z) \;\Rightarrow\; (x\, p\, z) &&\text{(bắc cầu)}\\
  &p \equiv q^{-1} \wedge (x\, p\, y) \;\Rightarrow\; (y\, q\, x) &&\text{(nghịch đảo)}
\end{align}
Năm luật này là tập con của ngữ nghĩa RDFS và OWL 2 RL đủ cho hai KG của bài. Luật \texttt{rdfs:domain} và
\texttt{rdfs:range} (suy ra kiểu của chủ thể từ vị từ) cố ý không bật: ở đây domain và range được dùng như
\emph{ràng buộc} để chặn dữ liệu sai, không phải để suy ra kiểu mới.

\maDai{kg_base__rdfs_closure}{Suy diễn tiến năm luật RDFS/OWL tới điểm bất động}

Trên đồ thị y tế, suy diễn thêm \SL{kg/health/inferred} bộ ba vào \SL{kg/health/rdf_triples} bộ ba gốc; trên
đồ thị tài chính thêm \SL{kg/finance/inferred} bộ ba vào \SL{kg/finance/rdf_triples}. Hình~\ref{fig:suydien}
minh họa trên một phần của \texttt{VietnamFinanceKG}: vị từ \textsf{THUOC\_VE} bắc cầu nên từ
``FPT thuộc VN30'' và ``VN30 thuộc HOSE'' suy ra ``FPT thuộc HOSE''.

\hinhtikz{t2_inference}{Suy diễn tiến trên \texttt{VietnamFinanceKG}: cạnh tường minh nét liền, bộ ba suy ra nét
đứt đỏ}{suydien}

\subsection{Truy vấn SPARQL}
SPARQL là ngôn ngữ truy vấn chuẩn của RDF~\cite{w3c2013sparql}: một truy vấn là một \emph{mẫu đồ thị} chứa biến
(\texttt{?thuoc}), và kết quả là mọi cách gán biến làm mẫu khớp với đồ thị. Bài chạy \SL{dem/sparql} truy vấn
trên cả đồ thị gốc lẫn đồ thị đã suy diễn (Bảng~\ref{tab:sparql}). Truy vấn Q1 đọc mức độ tương tác qua khối
reification; Q5 dùng đường thuộc tính \texttt{(p|q)+} để tìm mọi biến chịu tác động từ lãi suất Fed sau bất kỳ số
bước nào.

\maTruyVan{1}{Q1: thuốc tương tác nguy hiểm với rượu bia, đọc mức độ và hậu quả qua reification}
\maTruyVan{5}{Q5: mọi biến chịu tác động trực tiếp hay gián tiếp của lãi suất Fed (đường thuộc tính)}
\maTruyVan{2}{Q2: món ăn và chỉ số GI, truy vấn theo lớp cha \texttt{ex:Food}}

\bang{tab_sparql}{Kết quả \SL{dem/sparql} truy vấn SPARQL trước và sau suy diễn}{sparql}

Cột ``trước suy diễn'' cho thấy vai trò thực sự của bản thể học. Q2, Q4 và Q6 hỏi theo lớp cha (\textsf{Food},
\textsf{Drug}, \textsf{Stock}) trong khi dữ liệu chỉ ghi lớp lá (\textsf{ModerateFood}, \textsf{Biguanide}\dots),
nên trên đồ thị gốc chúng trả về rỗng; sau suy diễn kế thừa kiểu chúng trả về đầy đủ. Không có suy diễn, người
viết truy vấn phải tự liệt kê mọi lớp con, và truy vấn hỏng mỗi khi bản thể học thêm một lớp mới.
Hình~\ref{fig:kg_inference} tổng hợp hai hiệu ứng: số bộ ba tăng thêm và số dòng kết quả của từng truy vấn.

\hinhpdf{kg_inference}{1.0}{Ảnh hưởng của suy diễn: số bộ ba RDF trước, sau suy diễn (trái) và số dòng kết quả
của từng truy vấn SPARQL (phải)}{kg_inference}

\subsection{Kiểm tra ràng buộc}\label{sec:rangbuoc}
Để chứng minh ràng buộc domain/range thực sự hoạt động, kịch bản \texttt{build\_kgs.py} cố tình thêm
\SL{dem/violations} bộ ba sai. Cả bốn đều bị từ chối với thông báo nêu đúng lớp mong đợi và lớp thực tế
(Bảng~\ref{tab:vipham}). Đây là lớp phòng thủ đầu tiên chống lại dữ liệu bẩn, đặt ngay ở thời điểm nạp, trước
khi một bộ ba sai kịp lan vào câu trả lời của tác tử.

\bang{tab_violations}{Bốn bộ ba cố tình sai và thông báo của bộ kiểm tra ràng buộc}{vipham}

\section{Xây KB và KG từ dữ liệu bảng}
Không phải tri thức nào cũng có sẵn dạng văn bản. Bộ dữ liệu chẩn đoán đái tháo đường~\cite{mustafa2023diabetes}
là một bảng CSV \SL{kg/tabular/diabetes/records} dòng; dữ liệu giá cổ phiếu cũng là bảng. Tiếp cận của bài
tham khảo dungvt, đưa một bảng thành hai dạng tri thức song song (Hình~\ref{fig:kbpipeline}):

\begin{itemize}
  \item \textbf{KB văn bản}: mỗi dòng thành một câu mô tả, ví dụ ``Hồ sơ Person\_0: giới tính nữ, tuổi 80, tăng
        huyết áp không, bệnh tim có\dots''. Dạng này đưa vào RAG được.
  \item \textbf{KG bộ ba}: mỗi ô thành một bộ ba $(\textsf{Person\_0}, \textsf{has\_age}, 80)$. Dạng này truy vấn
        chính xác theo điều kiện số được.
\end{itemize}

\hinhtikz{t2_kb_pipeline}{Xây KB và KG từ dữ liệu bảng: mỗi dòng CSV được văn bản hóa thành một tài liệu (KB) và
tách thành các bộ ba (KG), phục vụ hai kiểu truy vấn khác nhau}{kbpipeline}

\maNguon{kg_tabular__diabetes_documents}{Văn bản hóa mỗi hồ sơ thành một câu tiếng Việt}
\maNguon{kg_tabular__diabetes_triples}{Tách mỗi hồ sơ thành các bộ ba thuộc tính}

\bang{tab_tabular}{Chuyển bảng chẩn đoán đái tháo đường thành KB văn bản và KG}{tabular}

Hai dạng tri thức trả lời hai loại câu hỏi khác nhau. Câu hỏi ``có bao nhiêu người trên 70 tuổi, tăng huyết áp,
HbA1c từ 8,8 trở lên'' là một phép lọc trên thuộc tính số: KG lọc ba điều kiện, trả đúng
\SL{kg/tabular/numeric_query/kg_hits} trên \SL{kg/tabular/numeric_query/truth} hồ sơ trong
\SL{kg/tabular/numeric_query/kg_ms} ms. Tìm kiếm văn bản bằng TF-IDF trên 20.000 tài liệu mẫu chỉ đưa được
\SL{kg/tabular/numeric_query/tfidf_top20_correct} hồ sơ đúng vào top-20 (Hình~\ref{fig:kg_tabular}), vì độ tương
đồng từ vựng không hiểu phép so sánh ``$\ge 8{,}8$''. Bài học: RAG mạnh với câu hỏi diễn giải, KG mạnh với câu
hỏi có điều kiện và số liệu; một hệ thống tốt cần cả hai.

\hinhpdf{kg_tabular}{1.0}{Câu hỏi có điều kiện số trên dữ liệu bảng: truy vấn KG trả đúng toàn bộ, tìm văn bản
TF-IDF gần như bỏ sót}{kg_tabular}

\section{So sánh chi tiết: Graph Database và Vector Database}
Hai loại kho dữ liệu chuyên dụng cho hệ thống dùng LLM là cơ sở dữ liệu đồ thị (Neo4j, GraphDB, Amazon
Neptune) và cơ sở dữ liệu vector (FAISS, Milvus, Qdrant, Chroma). Chúng trả lời hai câu hỏi khác nhau
(Hình~\ref{fig:graphvector}): Graph DB trả lời \emph{``cái gì nối với cái gì, bằng quan hệ gì''}; Vector DB trả
lời \emph{``cái gì giống cái này nhất''}.

\hinhtikz{t2_graph_vector}{Graph DB duyệt quan hệ tường minh giữa các thực thể; Vector DB tìm láng giềng gần nhất
theo khoảng cách trong không gian nhúng nhiều chiều}{graphvector}

\begin{table}[H]\centering
\caption{Bảng đối chuẩn kỹ thuật Graph Database và Vector Database}\label{tab:graphvector}
{\small\begin{tabular}{L{3.2cm}L{5.4cm}L{5.4cm}}
\toprule
Tiêu chí & Graph Database & Vector Database \\
\midrule
Đơn vị dữ liệu & nút, cạnh có nhãn, thuộc tính & vector dày $\mathbb{R}^d$ kèm metadata \\
Mô hình truy vấn & khớp mẫu đồ thị (Cypher, SPARQL), duyệt đường & láng giềng gần nhất (k-NN, ANN) \\
Chỉ mục & danh sách kề, chỉ mục theo nhãn & HNSW, IVF, lượng tử hóa PQ \\
Câu hỏi đa bước & tự nhiên: mỗi bước là một bước duyệt & không có khái niệm bước \\
Câu hỏi ngữ nghĩa mờ & kém: cần đúng tên thực thể & mạnh: hiểu diễn đạt khác nhau \\
Điều kiện số, phép đếm & chính xác & không hỗ trợ trực tiếp \\
Khả năng giải thích & đường đi và nguồn của từng cạnh & điểm tương đồng, khó giải thích \\
Chi phí xây dựng & cao: cần trích xuất, chuẩn hóa thực thể & thấp: nhúng văn bản thô \\
Cập nhật & thêm, sửa bộ ba cục bộ & nhúng lại đoạn văn đổi \\
Vai trò trong bài & \texttt{VietnamDiabetesKG}, \texttt{VietnamFinanceKG} & chỉ mục e5 trên \SL{retrieval_grid/n_docs} tài liệu \\
\bottomrule
\end{tabular}}
\end{table}

Hai loại kho không cạnh tranh mà bổ sung: Vector DB tìm \emph{đoạn văn} giải thích, Graph DB cung cấp
\emph{quan hệ} chính xác và đường suy luận. GraphRAG ở Chương~\ref{chap:rag} và tác tử ở Chương~\ref{chap:agent}
dùng cả hai cùng lúc.
